\documentclass[../../main.tex]{subfiles}
\begin{document}
\section{体积测量的几何}\label{section:ch16-the-geometry-of-volume-measurement}

\begin{figure}[ht]
\centering
\includegraphics[width=0.85\textwidth]{fig-16-1.png}
\caption{作为“带符号长度计”的余向量场}\label{figure:ch16-covector-length-meter}
\end{figure}

我们如何才能对多重积分作出与坐标无关的理解？首先注意，要在一个流形上以与坐标无关的方式积分取值于实数的\emph{函数}是不可能的，至少在不附加诸如 Riemann 度量之类额外结构的情况下如此。很容易看出为何如此，即使在最简单的情形中亦然：设 $C\subseteq\mathbb{R}^n$ 是一个闭球，$f:C\to\mathbb{R}$ 是恒等于 $1$ 的常值函数 $f(x)\equiv1$。则
\begin{align*}
\int_C f\,\mathrm{d}V=\operatorname{Vol}(C),
\end{align*}
它显然在坐标变换下不是不变的，哪怕我们只限于线性坐标变换。

另一方面，在\cref{chapter:ch11}中我们已经说明，余向量场可以沿曲线以自然的方式被积分。让我们更几何地思考一下为何如此。流形 $M$ 上的一个余向量场对每个切向量指定一个数，其指定方式是：把切向量乘以一个常数，就会使所得的数乘以同一个常数。因此，余向量场可以设想为给每个切空间的每个一维子空间指定一把“带符号的长度计”（\cref{figure:ch16-covector-length-meter}），而且这种指定与坐标无关。计算一个余向量场的线积分，实际上就是沿着曲线各点使用这把随点变化的尺度，从而给一条曲线指定一个“长度”。

现在我们希望寻找一种“场”，它能够以与坐标无关的方式在维数 $k>1$ 的子流形上被积分。它在每一点处的值应当能够被解释为切空间的某个 $k$ 维子空间上的一把“带符号体积计”：一台机器 $\omega$，它在一点处接受任意 $k$ 个切向量 $(v_1,\dots,v_k)$，并返回一个数 $\omega(v_1,\dots,v_k)$，我们可以把这个数想作由这些向量张成的平行多面体的“带符号体积”，其度量依据由 $\omega$ 所确定的尺度进行。

这种机器最显而易见的例子就是 $\mathbb{R}^n$ 中的行列式。例如，大多数线性代数教材都证明了，对于任意两个向量 $v_1,v_2\in\mathbb{R}^2$，$\det(v_1,v_2)$ 至多差一个符号就是由 $v_1,v_2$ 张成的平行四边形的面积。不难证明（见习题~\cref{exercise:ch16-01}），类似的事实在所有维数下都成立。请记住，行列式正是一个斜对称张量的例子。

让我们来考虑一下，我们对一台一般的“带符号 $k$ 维体积计” $\omega$ 应当期望它具有哪些性质。为了与我们关于体积的直觉相一致，用一个常数去乘任意一个向量，应当使体积乘以同一个常数（\cref{figure:ch16-scaling-by-constant}）；而把 $k$ 维平行多面体在第 $i$ 个位置上加入两个向量之和所得到的体积，应当等于原来在第 $i$ 个位置上的两个平行多面体的体积之和（\cref{figure:ch16-sum-of-two-vectors}）：
\begin{align}
\omega(v_1,\dots,c\,v_i,\dots,v_k)&=c\,\omega(v_1,\dots,v_i,\dots,v_k),\\
\omega(v_1,\dots,v_i+v_i',\dots,v_k)&=\omega(v_1,\dots,v_i,\dots,v_k)+\omega(v_1,\dots,v_i',\dots,v_k).
\end{align}
（注意\cref{figure:ch16-sum-of-two-vectors} 中的向量都假定位于同一平面内。）这表明 $\omega$ 应当是多重线性的，因而应当是一个协变 $k$-张量。

还有一条我们应当期望的性质：既然一个线性相关的 $k$-元组张成的平行多面体具有零的 $k$ 维体积，那么当 $\omega$ 作用于这样一个 $k$-元组时就应当给出零。由\cref{lemma:ch14-alternating-equiv}，这就迫使 $\omega$ 是斜对称的。于是，斜对称张量场就成为可与坐标无关地进行积分的有希望的对象。在本章中，我们将说明如何做到这一点。

\begin{figure}[ht]
\centering
\begin{minipage}[t]{0.48\textwidth}
\centering
\includegraphics[width=0.9\textwidth]{fig-16-2.png}
\caption{乘以一个常数}\label{figure:ch16-scaling-by-constant}
\end{minipage}
\end{figure}

\begin{figure}[ht]
\centering
\begin{minipage}[t]{0.48\textwidth}
\centering
\includegraphics[width=0.9\textwidth]{fig-16-3.png}
\caption{两个向量之和}\label{figure:ch16-sum-of-two-vectors}
\end{minipage}
\end{figure}
\end{document}
